% GENERATED FILE. Edit canonical lessons, then run npm run generate:books.
% canonical-source-hash: fnv1a64:2a540769a7a05b21
% canonical-lessons: UR-C330-safha, UR-C330-unvan, UR-C330-qari, UR-C330-bikna, UR-C330-tarjuma, UR-R330-review-324-330, UR-R330-review-final-330

\chapter{Turning the Pages}
\label{ch:ur-turning-the-pages}
\hlchaptermodality{330}

\begin{chapteropening}
\textbf{By the end of this chapter:} \emph{I can describe a book and say how well it sold.}

Complete the last lesson of chapter 330: I can describe a book and say how well it sold.
\end{chapteropening}

\section[\texorpdfstring{\ur{صفحہ} (safhā)}{صفحہ (safhā)}]{\ur{صفحہ} (safhā) — a page}

\label{lesson:UR-C330-safha}



\begin{quote}
Before the new one: say the Urdu for a chapter of a book, then the Urdu for an ending.
\end{quote}

\subsection*{You'll want to know: \ur{صفحہ}}
\textbf{\ur{صفحہ}} — \emph{safhā} — \textquotedblleft{}a page\textquotedblright{}.

\begin{quote}
\ur{اس} \ur{کتاب} \ur{میں} \ur{سو} \ur{صفحے} \ur{ہیں۔}
\end{quote}

\begin{quote}
\emph{is kitāb meṅ sau safhe haiṅ.}
\end{quote}

\begin{quote}
This book has a hundred pages.
\end{quote}

Masculine.

\begin{cousinweb}
\textbf{\ur{صفحہ}} (\emph{safhā}) is usually traced, through Persian, to Arabic \emph{ṣafḥa}, \textquotedblleft{}a flat surface, a page\textquotedblright{}. The plural is \textbf{\ur{صفحے}}.
\end{cousinweb}

\begin{grammarlens}[title={What you've built}]
One more word for telling what happened.
\end{grammarlens}

\subsection*{Your turn}
\begin{itemize}
\raggedright
  \item \emph{Say it:} \emph{safhā}
  \item \emph{Say it:} \emph{safhā}, once more
  \item \emph{Say it:} \emph{is kitāb meṅ sau safhe haiṅ.} — This book has a hundred pages.
  \item \emph{Recall:} say the Urdu for a chapter of a book, then the Urdu for an ending, then say \emph{safhā} again
\end{itemize}

\begin{tcolorbox}[breakable,colback=teal!4,colframe=teal!35!black,title={Before you move on}]
What does \ur{صفحہ} mean? (\textquotedblleft{}A page\textquotedblright{}.) Say it once more.
\end{tcolorbox}
\section[\texorpdfstring{\ur{عنوان} (unvān)}{عنوان (unvān)}]{\ur{عنوان} (unvān) — a title of a book; a heading}

\label{lesson:UR-C330-unvan}



\begin{quote}
Before the new one: say the Urdu for an ending, then the Urdu for a page.
\end{quote}

\subsection*{You'll want to know: \ur{عنوان}}
\textbf{\ur{عنوان}} — \emph{unvān} — \textquotedblleft{}a title of a book; a heading\textquotedblright{}.

\begin{quote}
\ur{کتاب} \ur{کا} \ur{عنوان} \ur{کیا} \ur{ہے؟}
\end{quote}

\begin{quote}
\emph{kitāb kā unvān kyā hai?}
\end{quote}

\begin{quote}
What is the book's title?
\end{quote}

Masculine.

\begin{cousinweb}
\textbf{\ur{عنوان}} (\emph{unvān}) is usually traced, through Persian, to Arabic \emph{ʿunwān}, \textquotedblleft{}the heading of a letter, a title\textquotedblright{}.
\end{cousinweb}

\begin{grammarlens}[title={What you've built}]
One more word for telling what happened.
\end{grammarlens}

\subsection*{Your turn}
\begin{itemize}
\raggedright
  \item \emph{Say it:} \emph{unvān}
  \item \emph{Say it:} \emph{unvān}, once more
  \item \emph{Say it:} \emph{kitāb kā unvān kyā hai?} — What is the book's title?
  \item \emph{Recall:} say the Urdu for an ending, then the Urdu for a page, then say \emph{unvān} again
\end{itemize}

\begin{tcolorbox}[breakable,colback=teal!4,colframe=teal!35!black,title={Before you move on}]
What does \ur{عنوان} mean? (\textquotedblleft{}A title of a book; a heading\textquotedblright{}.) Say it once more.
\end{tcolorbox}
\section[\texorpdfstring{\ur{قاری} (qārī)}{قاری (qārī)}]{\ur{قاری} (qārī) — a reader}

\label{lesson:UR-C330-qari}



\begin{quote}
Before the new one: say the Urdu for a page, then the Urdu for a title of a book.
\end{quote}

\subsection*{You'll want to know: \ur{قاری}}
\textbf{\ur{قاری}} — \emph{qārī} — \textquotedblleft{}a reader\textquotedblright{}.

\begin{quote}
\ur{یہ} \ur{کتاب} \ur{ہر} \ur{قاری} \ur{کو} \ur{پسند} \ur{آئی۔}
\end{quote}

\begin{quote}
\emph{yih kitāb har qārī ko pasand āī.}
\end{quote}

\begin{quote}
Every reader liked this book.
\end{quote}

Masculine or feminine, by the person meant.

\begin{cousinweb}
\textbf{\ur{قاری}} (\emph{qārī}) is usually traced, through Persian, to Arabic \emph{qāriʾ}, \textquotedblleft{}one who reads aloud\textquotedblright{}; in Urdu it is any reader of a book or a paper.
\end{cousinweb}

\begin{grammarlens}[title={What you've built}]
One more word for telling what happened.
\end{grammarlens}

\subsection*{Your turn}
\begin{itemize}
\raggedright
  \item \emph{Say it:} \emph{qārī}
  \item \emph{Say it:} \emph{qārī}, once more
  \item \emph{Say it:} \emph{yih kitāb har qārī ko pasand āī.} — Every reader liked this book.
  \item \emph{Recall:} say the Urdu for a page, then the Urdu for a title of a book, then say \emph{qārī} again
\end{itemize}

\begin{tcolorbox}[breakable,colback=teal!4,colframe=teal!35!black,title={Before you move on}]
What does \ur{قاری} mean? (\textquotedblleft{}A reader\textquotedblright{}.) Say it once more.
\end{tcolorbox}
\section[\texorpdfstring{\ur{بکنا} (biknā)}{بکنا (biknā)}]{\ur{بکنا} (biknā) — to be sold; to sell}

\label{lesson:UR-C330-bikna}



\begin{quote}
Before the new one: say the Urdu for a title of a book, then the Urdu for a reader.
\end{quote}

\subsection*{You'll want to know: \ur{بکنا}}
\textbf{\ur{بکنا}} — \emph{biknā} — \textquotedblleft{}to be sold; to sell\textquotedblright{}.

\begin{quote}
\ur{اس} \ur{کی} \ur{کتاب} \ur{جلدی} \ur{بک} \ur{گئی۔}
\end{quote}

\begin{quote}
\emph{us kī kitāb jaldī bik gaī.}
\end{quote}

\begin{quote}
His book sold quickly.
\end{quote}

Past (it …): \ur{وہ} \ur{کتاب} \ur{بکی}, \emph{voh kitāb bikī}.

\begin{cousinweb}
is inherited: CDIAL lists the Hindi cousin \emph{biknā}, \textquotedblleft{}to be sold\textquotedblright{}, under Sanskrit \emph{vikrīyatē} (CDIAL 11642). It is the partner of \textbf{\ur{بیچنا}}: you \textbf{\ur{بیچنا}} a thing, and it \textbf{\ur{بکنا}}.
\end{cousinweb}

\begin{grammarlens}[title={What you've built}]
One more word for telling what happened.
\end{grammarlens}

\subsection*{Your turn}
\begin{itemize}
\raggedright
  \item \emph{Say it:} \emph{biknā}
  \item \emph{Say it:} \emph{biknā}, once more
  \item \emph{Say it:} \emph{us kī kitāb jaldī bik gaī.} — His book sold quickly.
  \item \emph{Recall:} say the Urdu for a title of a book, then the Urdu for a reader, then say \emph{biknā} again
\end{itemize}

\begin{tcolorbox}[breakable,colback=teal!4,colframe=teal!35!black,title={Before you move on}]
What does \ur{بکنا} mean? (\textquotedblleft{}To be sold; to sell\textquotedblright{}.) Say it once more.
\end{tcolorbox}
\section[\texorpdfstring{\ur{ترجمہ} (tarjumā)}{ترجمہ (tarjumā)}]{\ur{ترجمہ} (tarjumā) — a translation}

\label{lesson:UR-C330-tarjuma}



\begin{quote}
Before the new one: say the Urdu for a reader, then the Urdu for to be sold.
\end{quote}

\subsection*{You'll want to know: \ur{ترجمہ}}
\textbf{\ur{ترجمہ}} — \emph{tarjumā} — \textquotedblleft{}a translation\textquotedblright{}.

\begin{quote}
\ur{اس} \ur{ناول} \ur{کا} \ur{انگریزی} \ur{میں} \ur{ترجمہ} \ur{ہوا۔}
\end{quote}

\begin{quote}
\emph{is nāvil kā angrezī meṅ tarjumā huā.}
\end{quote}

\begin{quote}
This novel was translated into English.
\end{quote}

Masculine.

\begin{cousinweb}
\textbf{\ur{ترجمہ}} (\emph{tarjumā}) is usually traced, through Persian, to Arabic \emph{tarjama}, \textquotedblleft{}a rendering into another tongue\textquotedblright{}. \textbf{\ur{ترجمہ} \ur{کرنا}} is to translate.
\end{cousinweb}

\begin{grammarlens}[title={What you've built}]
One more word for telling what happened.
\end{grammarlens}

\subsection*{Your turn}
\begin{itemize}
\raggedright
  \item \emph{Say it:} \emph{tarjumā}
  \item \emph{Say it:} \emph{tarjumā}, once more
  \item \emph{Say it:} \emph{is nāvil kā angrezī meṅ tarjumā huā.} — This novel was translated into English.
  \item \emph{Recall:} say the Urdu for a reader, then the Urdu for to be sold, then say \emph{tarjumā} again
\end{itemize}

\begin{tcolorbox}[breakable,colback=teal!4,colframe=teal!35!black,title={Before you move on}]
What does \ur{ترجمہ} mean? (\textquotedblleft{}A translation\textquotedblright{}.) Say it once more.
\end{tcolorbox}
\section[(dialogue)]{Looking back: from A Tale of Long Ago to Turning the Pages}

\label{lesson:UR-R330-review-324-330}



\begin{quote}
Say the Urdu for to tell a story and an epic tale.
\end{quote}

\subsection*{Your turn}
Say these aloud:

\begin{itemize}
\raggedright
  \item to tell a story, an epic tale, an anecdote, a riddle, a proverb
  \item a giant in tales, a fairy, magic, a ghost, a witch
  \item a curse, to rail at, to regret, to sulk, torment
  \item a treasure, to bury, a cobra, to bite (of a snake), a diamond
  \item a prince, a disguise, to put to the test, deceit, to lead astray
  \item a short story, an author, a character (in a story), a chapter of a book, an ending
  \item a page, a title of a book, a reader, to be sold, a translation
\end{itemize}

\begin{tcolorbox}[breakable,colback=teal!4,colframe=teal!35!black,title={Before you move on}]
To tell a story? (\textbf{\ur{سنانا}}, \emph{sunānā}.) An epic tale? (\textbf{\ur{داستان}}, \emph{dāstān}.) An anecdote? (\textbf{\ur{قصہ}}, \emph{qissā}.) A riddle? (\textbf{\ur{پہیلی}}, \emph{pahelī}.) A proverb? (\textbf{\ur{کہاوت}}, \emph{kahāvat}.) A giant in tales? (\textbf{\ur{دیو}}, \emph{dev}.) A fairy? (\textbf{\ur{پری}}, \emph{parī}.) Magic? (\textbf{\ur{جادو}}, \emph{jādū}.) A ghost? (\textbf{\ur{بھوت}}, \emph{bhūt}.) A witch? (\textbf{\ur{ڈائن}}, \emph{ḍāin}.) A curse? (\textbf{\ur{لعنت}}, \emph{lānat}.) To rail at? (\textbf{\ur{کوسنا}}, \emph{kosnā}.) To regret? (\textbf{\ur{پچھتانا}}, \emph{pachhtānā}.) To sulk? (\textbf{\ur{روٹھنا}}, \emph{rūṭhnā}.) Torment? (\textbf{\ur{عذاب}}, \emph{azāb}.) A treasure? (\textbf{\ur{خزانہ}}, \emph{khazānā}.) To bury? (\textbf{\ur{گاڑنا}}, \emph{gāṛnā}.) A cobra? (\textbf{\ur{ناگ}}, \emph{nāg}.) To bite (of a snake)? (\textbf{\ur{ڈسنا}}, \emph{ḍasnā}.) A diamond? (\textbf{\ur{ہیرا}}, \emph{hīrā}.) A prince? (\textbf{\ur{شہزادہ}}, \emph{shahzādā}.) A disguise? (\textbf{\ur{بھیس}}, \emph{bhes}.) To put to the test? (\textbf{\ur{آزمانا}}, \emph{āzmānā}.) Deceit? (\textbf{\ur{فریب}}, \emph{fareb}.) To lead astray? (\textbf{\ur{بہکانا}}, \emph{bahkānā}.) A short story? (\textbf{\ur{افسانہ}}, \emph{afsānā}.) An author? (\textbf{\ur{مصنف}}, \emph{musannif}.) A character (in a story)? (\textbf{\ur{کردار}}, \emph{kirdār}.) A chapter of a book? (\textbf{\ur{باب}}, \emph{bāb}.) An ending? (\textbf{\ur{انجام}}, \emph{anjām}.) A page? (\textbf{\ur{صفحہ}}, \emph{safhā}.) A title of a book? (\textbf{\ur{عنوان}}, \emph{unvān}.) A reader? (\textbf{\ur{قاری}}, \emph{qārī}.) To be sold? (\textbf{\ur{بکنا}}, \emph{biknā}.) A translation? (\textbf{\ur{ترجمہ}}, \emph{tarjumā}.)
\end{tcolorbox}
\section[(dialogue)]{Turning the Pages — once more}

\label{lesson:UR-R330-review-final-330}



\begin{quote}
Say the Urdu for a page and a title of a book.
\end{quote}

\subsection*{Your turn}
\begin{itemize}
\raggedright
  \item \emph{Say it:} a page, a title of a book, a reader, to be sold, a translation
\end{itemize}

\begin{tcolorbox}[breakable,colback=teal!4,colframe=teal!35!black,title={Before you move on}]
A page? (\textbf{\ur{صفحہ}}, \emph{safhā}.) A title of a book? (\textbf{\ur{عنوان}}, \emph{unvān}.) A reader? (\textbf{\ur{قاری}}, \emph{qārī}.) To be sold? (\textbf{\ur{بکنا}}, \emph{biknā}.) A translation? (\textbf{\ur{ترجمہ}}, \emph{tarjumā}.)
\end{tcolorbox}
